\documentclass[11pt,a4paper]{article}

\usepackage[utf8]{inputenc}
\usepackage[T1]{fontenc}
\usepackage[spanish,es-noshorthands]{babel}
\usepackage{xcolor}
\usepackage[most]{tcolorbox}
\usepackage{listings}
\usepackage{fancyhdr}
\usepackage{lastpage}
\usepackage{enumitem}
\usepackage[hidelinks]{hyperref}

\setlength{\headheight}{30pt}
\pagestyle{fancy}
\fancyhf{}
\renewcommand{\headrulewidth}{0pt}
\fancyhead[R]{%
  \fontsize{8}{9.6}\selectfont
  Licenciatura en Ciencias de la Computación\\
  Sistemas Operativos\\[2pt]
  \rule{4.2cm}{0.5pt}
}
\fancyfoot[C]{\thepage\ / \pageref{LastPage}}

\newtcolorbox{trackbox}[1]{
  colback=green!7!white, colframe=green!45!black,
  arc=2mm, boxrule=0.6pt, left=8pt, right=8pt, top=6pt, bottom=6pt,
  title=TRACK -- #1, fonttitle=\bfseries, coltitle=green!35!black,
  colbacktitle=green!7!white, boxsep=1pt, breakable
}

\newtcolorbox{bonusbox}{
  colback=red!7!white, colframe=red!55!black,
  arc=2mm, boxrule=0.6pt, left=8pt, right=8pt, top=6pt, bottom=6pt,
  title=BONUS -- POR PUNTOS, fonttitle=\bfseries, coltitle=red!55!black,
  colbacktitle=red!7!white, boxsep=1pt, breakable
}

\newtcolorbox{extrabox}{
  colback=orange!8!white, colframe=orange!55!white,
  arc=2mm, boxrule=0.6pt, left=8pt, right=8pt, top=6pt, bottom=6pt,
  title=EXTRA, fonttitle=\bfseries, coltitle=orange!55!black,
  colbacktitle=orange!8!white, boxsep=1pt, breakable
}

\lstdefinestyle{codestyle}{
  backgroundcolor=\color{black!5}, rulecolor=\color{black!25},
  frame=single, framerule=0.5pt, basicstyle=\ttfamily\small,
  breaklines=true, showstringspaces=false, tabsize=4,
  keywordstyle=\color{blue!60!black}\bfseries,
  commentstyle=\color{green!40!black},
  stringstyle=\color{orange!70!black},
  columns=fullflexible, upquote=true
}
\lstset{style=codestyle}
\setlist{nosep}

\title{Trabajo Práctico 5 -- Ejercicios extra}
\author{}
\date{}

\begin{document}
\renewcommand{\contentsname}{Tracks}
\maketitle
\thispagestyle{fancy}
\begin{center}
SISTEMAS OPERATIVOS
\end{center}
\newpage

\tableofcontents
\newpage

\section{Límite de instancias con un semáforo (7/10)}

\begin{trackbox}{7/10}
\textbf{Descripción:} limitar la cantidad de instancias concurrentes de un
programa usando un semáforo contador de System V.
\end{trackbox}

Implementá un programa que limite la cantidad de instancias concurrentes usando
un semáforo de System V. Interfaz sugerida:

\begin{lstlisting}[language=bash]
instance-limit SEM_KEY MAX_INSTANCES -- COMMAND [ARGS...]
\end{lstlisting}

Investigá la inicialización, la reserva y liberación de cupos, el
comportamiento cuando se alcanza el límite y la limpieza del semáforo.

\newpage
\section{Aproximación de $\pi$ y rendimiento (5/10)}

\begin{trackbox}{5/10}
\textbf{Descripción:} aproximar $\pi$ sumando términos de una expansión y
medir el costo de aumentar la cantidad de términos.
\end{trackbox}

Implementá dos versiones de un programa que reciba $N$ y calcule:

\[
\pi_N = 4 \sum_{k=0}^{N-1} \frac{(-1)^k}{2k+1}.
\]

Una versión debe ser secuencial. La otra debe repartir los términos entre
varios trabajadores concurrentes, recibir una cantidad configurable de
trabajadores y combinar los resultados parciales.

El error de esta aproximación cumple

\[
\left|\pi - \pi_N\right| < \frac{4}{2N+1}.
\]

Mostrá la aproximación, el error y el tiempo de ejecución de ambas versiones.
Probá distintos valores de $N$ y distintas cantidades de trabajadores;
compará los resultados y el costo de la coordinación.

\begin{bonusbox}
Conservá las versiones secuencial y multi-worker de la consigna y agregá una
tercera versión usando OpenMP. Las tres deben calcular la misma aproximación,
aceptar parámetros equivalentes y producir resultados comparables. Compará
las tres ejecuciones con distintas cantidades de trabajadores o hilos e
informá el speedup. Se sugiere usar C o C++.
\end{bonusbox}

\newpage
\section{Transferencia por pipes y SSH (4/10)}

\begin{trackbox}{4/10}
\textbf{Descripción:} transferir y extraer un directorio remoto usando un
pipe, sin crear un archivo intermedio.
\end{trackbox}

Investigá brevemente qué resuelven \texttt{scp} y \texttt{rsync}, y compará
esas herramientas con una transferencia basada en un flujo de datos.

Después, escribí un \texttt{one-liner} que:

\begin{itemize}
  \item empaquete un directorio local y envíe el resultado por la salida
  estándar;
  \item use \texttt{ssh} para transportar ese flujo;
  \item extraiga el flujo en un directorio de la máquina remota.
\end{itemize}

La solución debe usar pipes, \texttt{ssh} y \texttt{tar}. No uses
\texttt{scp} ni \texttt{rsync} para resolverla, y no generes un archivo
intermedio. Probá la línea con una máquina remota y verificá que la estructura
y el contenido lleguen correctamente.

\newpage
\section{Cliente-servidor: solicitud y respuesta (6/10)}

\begin{trackbox}{6/10}
\textbf{Descripción:} implementar una comunicación básica entre un proceso
cliente y un proceso servidor.
\end{trackbox}

Implementá dos procesos:

\begin{itemize}
  \item el servidor espera una solicitud, la procesa y envía una respuesta;
  \item el cliente envía la solicitud y muestra la respuesta recibida.
\end{itemize}

Elegí una operación sencilla, por ejemplo, solicitar al servidor que tire un
dado o enviarle una cuenta aritmética para que devuelva el resultado. Definí
un protocolo simple y asegurá al menos un intercambio completo de solicitud y
respuesta mediante un mecanismo de comunicación entre procesos.

\begin{bonusbox}
Convertí la solución en un juego mínimo con dos procesos. El cliente debe
mostrar un rectángulo grande usando secuencias de escape ANSI y un personaje
dentro del área. Configurá la entrada del terminal sin modo canónico ni eco
para leer las teclas \texttt{W}, \texttt{A}, \texttt{S} y \texttt{D}
inmediatamente, y restaurá la configuración al salir.
Como referencia para las secuencias de control del terminal, consultá
\href{https://www.xfree86.org/current/ctlseqs.html}{Xterm Control Sequences}.

Cada tecla debe enviarse al servidor. El servidor debe validar y aplicar el
movimiento, devolver el estado actualizado, incluyendo la posición del
personaje, y mantener la autoridad exclusiva sobre el estado. El cliente debe
renderizar en tiempo real únicamente el estado recibido.
\end{bonusbox}

\newpage
\section{Heartbeat y watchdog (6/10)}

\begin{trackbox}{6/10}
\textbf{Descripción:} detectar la ausencia de un proceso mediante mensajes
periódicos de heartbeat.
\end{trackbox}

Implementá dos procesos:

\begin{itemize}
  \item un proceso de trabajo que envíe periódicamente un mensaje de heartbeat;
  \item un watchdog que reciba esos mensajes y controle el tiempo transcurrido
  desde el último heartbeat.
\end{itemize}

Si el watchdog no recibe un heartbeat dentro de un timeout configurable, debe
informar que el proceso dejó de responder. Usá un mecanismo de comunicación
entre procesos, definí el formato de los mensajes y asegurá el cierre correcto
de los descriptores.

Probá tanto el funcionamiento normal como la ausencia o terminación del
proceso de trabajo.

\newpage
\section{Comunicación entre pares con FIFOs nombradas (5/10)}

\begin{trackbox}{5/10}
\textbf{Descripción:} construir una comunicación local entre procesos pares
usando FIFOs (\emph{First In, First Out}) nombradas.
\end{trackbox}

Construí una aplicación de mensajería local. Una sala se representa con una
carpeta compartida y cada proceso crea dentro de ella su propia FIFO nombrada.
Los procesos deben descubrirse y enviarse mensajes directamente, sin un
servidor central. Definí un protocolo simple, separá los mensajes sucesivos y
limpiá la FIFO al salir. Podés usar Protocol Buffers (protobuf) para
serializar los mensajes.

Probá varios pares en una misma sala, la entrada y salida de procesos y la
desaparición inesperada de uno de ellos. Observá las FIFOs y sus descriptores
con \texttt{ls -l} y \texttt{/proc/<pid>/fd}.

\begin{bonusbox}
Convertí la aplicación en una sala de chat local entre pares. Agregá una
interfaz de terminal cómoda que permita escribir mensajes y mostrar los
mensajes entrantes, manteniendo la comunicación directa, las salas por
carpetas y una salida limpia. Podés consultar la documentación
\href{https://www.xfree86.org/current/ctlseqs.html}{Xterm Control Sequences}
para las secuencias de control del terminal.
\end{bonusbox}

\end{document}
